% Figure: a block resting on a block, the lower one pulled. Free-body
% diagrams of both blocks side by side, showing the friction that must
% accelerate the upper block and the limit at which it slips.
\begin{figure}[htbp]
\centering
\begin{tikzpicture}[
  force/.style={draw=engblue, -{Stealth}, thick},
  wt/.style={draw=engblack, -{Stealth}, thick},
  fr/.style={draw=engamber, -{Stealth}, thick},
  nrm/.style={draw=engred, -{Stealth}, thick},
  lbl/.style={font=\scriptsize},
]
% left: the physical picture
\begin{scope}
  \draw[draw=enggray, thick] (-0.3,0) -- (5,0);
  \fill[enggray!12] (-0.3,0) rectangle (5,-0.35);
  % lower block m1
  \draw[draw=engblack, thick, fill=engblue!8] (0.4,0) rectangle (3.4,1.0);
  \node[lbl] at (1.9,0.5) {$m_1$};
  % upper block m2
  \draw[draw=engblack, thick, fill=englime!18] (1.2,1.0) rectangle (2.9,1.8);
  \node[lbl] at (2.05,1.4) {$m_2$};
  % applied pull on lower block
  \draw[force] (3.4,0.5) -- ++(1.3,0) node[right, lbl, engblue] {$P$};
  \node[lbl, anchor=north] at (2.3,-0.15) {applied to lower block};
\end{scope}
% right: stacked free-body diagrams
\begin{scope}[shift={(7.6,0)}]
  % upper block FBD
  \draw[draw=engblack, thick, fill=englime!18] (-0.7,1.6) rectangle (0.7,2.4);
  \node[lbl] at (0,2.0) {$m_2$};
  \draw[wt] (0,1.6) -- ++(0,-0.9) node[below, lbl] {$m_2 g$};
  \draw[nrm] (0,2.4) -- ++(0,0.9) node[above, lbl, engred] {$N_2$};
  \draw[fr] (0,2.0) -- ++(1.1,0) node[right, lbl, engamber] {$f$ (drives $m_2$)};
  % lower block FBD
  \draw[draw=engblack, thick, fill=engblue!8] (-1.0,-0.7) rectangle (1.0,0.3);
  \node[lbl] at (-0.4,-0.2) {$m_1$};
  \draw[wt] (-0.4,-0.7) -- ++(0,-0.8) node[below, lbl] {$m_1 g$};
  \draw[nrm] (-0.4,0.3) -- ++(0,0.8) node[above, lbl, engred] {$N_1$};
  \draw[force] (1.0,-0.2) -- ++(1.1,0) node[right, lbl, engblue] {$P$};
  \draw[fr] (0.4,0.3) -- ++(-1.1,0) node[left, lbl, engamber] {$f$ (reaction)};
  \draw[nrm] (0.4,0.3) -- ++(0,-0.7) node[pos=1.0, below, lbl, engred] {$N_2$};
\end{scope}
\end{tikzpicture}
\caption[Two stacked blocks and their free-body diagrams]{A block
$m_2$ rides on a block $m_1$ that is pulled by a force $P$ along a
frictionless floor. The only horizontal force on the upper block is the
friction $f$ from the lower block's top face, so that friction is what
must accelerate $m_2$; it can grow no larger than $\mu_s m_2 g$. While
$P$ is small the two blocks share one acceleration and the friction
stays below its limit; when $P$ is large enough that the shared
acceleration would need more friction than $\mu_s m_2 g$, the upper
block slips and the two accelerations part. The paired diagrams keep
the third-law friction couple straight: the lower block feels the equal
and opposite reaction $f$ backward on its top face.}
\label{fig:vol03ch04:block-on-block}
\end{figure}
